\subsection{LENGTH OF AN INTERVAL}

DEFINITION 3. The length of a bounded interval with end-points \(a\) and \(b\) ( \(a \leqslant b\) ) is defined to be \(b - a\) .

Every bounded interval which contains more than one point therefore has length \(> 0\) . If \(a \leqslant b\) , the four intervals \([a, b]\) , \(]a, b[\) , \([a, b[\) and \(]a, b[\) all have the same length. An interval with end-points \(a + c\) and \(b + c\) has the same length as an interval with end-points \(a\) and \(b\) ; in other words, the length of an interval is invariant under translation.

If \(a \leqslant c \leqslant d \leqslant b\) we have \(d - c \leqslant b - a\) . Hence if a bounded interval \(I\) is contained in a bounded interval \(I'\) , the length of \(I\) is less than or equal to the length of \(I'\) .

If \(n\) mutually disjoint open intervals \(I_1, I_2, \ldots, I_n\) are contained in the interval \([a, b] (a < b)\) it is easily seen by induction on \(n\) that, if \(I_k = ]c_k, d_k[\) , there is a permutation \(\sigma\) of the indices \(k (1 \leqslant k \leqslant n)\) such that \(d_{\sigma(k)} \leqslant c_{\sigma(k+1)}\) for \(1 \leqslant k \leqslant n - 1\) . It follows immediately that the sum of the lengths of the intervals \(I_k\) is at most equal to the length of \([a, b]\) , and that equality holds if and only if \(c_{\sigma(1)} = a, d_{\sigma(n)} = b\) and \(d_{\sigma(k)} = c_{\sigma(k+1)}\) for \(1 \leqslant k \leqslant n - 1\) .

